\section*{Ethics Considerations}

\noindent
\textbf{Responsible disclosure.}
We reported the vulnerability to the Bluetooth Special Interest Group~(SIG) security team on
July~15, 2024, under the heading of a unilateral key downgrade and covering two scenarios, of
which the one this paper reports is the second.
On February~7, 2025 the SIG replied with the results of its internal analysis, considering the
scenario this paper reports to be practical.

\noindent
\textbf{Disclosure channel.}
UMD is a weakness in how the Bluetooth Classic specification permits key-length negotiation, not a
defect in one product, so we disclosed it to the SIG rather than to individual manufacturers, and
we did not contact any manufacturer directly.
Notification of affected members went through the SIG, which asked for a disclosure timeline on
their behalf and which we authorized to pass our plans on.

\noindent
\textbf{Experimental setup.}
All analysis in this paper is conducted entirely within the symbolic model using \Tamarin{}.
No real Bluetooth devices were involved in the discovery or validation of the UMD attack.
The attack traces are produced by the automated verifier and do not require physical proximity to
any device.
We have therefore not demonstrated the attack against a physical device, and the SIG's February
2025 reply invites exactly such a demonstration.

\noindent
\textbf{User-side risk assessment.}
The UMD attack requires a Bluetooth Classic peer that still negotiates the encryption key length
in the legacy manner and accepts a session key of eight octets or fewer. That is the same
precondition as the published key-negotiation downgrade attacks~\cite{antonioli2019knob}, so the
devices at risk are the old ones already at risk from those attacks, not a new population.
After the KNOB disclosure the SIG updated the Core Specification to ``recommend a minimum
encryption key length of 7 octets for BR/EDR connections''~\cite{sig2019knob}. That is a
recommendation rather than a requirement, and it does not exclude the key lengths this attack
needs.
Users can reduce exposure by: (i)~avoiding pairing with untrusted or unknown devices,
(ii)~preferring devices that enforce Secure Connections and maximum key lengths on both sides, and
(iii)~monitoring for unexpected re-pairing events that may indicate an active downgrade attack.

